% Source: https://artofproblemsolving.com/community/c3h1657695_2018_ms_mte_theme_round_all
% Stable identity: aops-topic-1657695-q12
Since $x$ is a root of $x^2+k$, $k=-x^2=-(a+b-2c)^2$, so we want $|a+b-2c|=(c-a)+(c-b)$ as small as possible. As the triangle is scalene with $c$ largest, $c-a$ and $c-b$ are different positive integers, so their sum is at least $1+2=3$, and the sum $3$ means the sides are $c-2,c-1,c$: consecutive, which is excluded. The sum cannot be $3$ otherwise, so it is at least $4$, attained by $\{c-a,c-b\}=\{1,3\}$, e.g.\ sides $2,4,5$ (a valid triangle as $2+4>5$).

\begin{asy}
usepackage("amsmath");
size(4.5cm);
defaultpen(fontsize(10pt));
real a=4,bb=2,c=5;
real cx=(bb*bb+c*c-a*a)/(2*c);
pair A=(0,0),B=(c,0),C=(cx,sqrt(bb*bb-cx*cx));
draw(A--B--C--cycle,linewidth(0.9));
label("$5$",(A+B)/2,S);label("$2$",(A+C)/2,W);label("$4$",(B+C)/2,NE);
label("\mbox{}",(min(currentpicture,user=true)+max(currentpicture,user=true))/2,black); // keeps the text after the figure black when LaTeX typesets the labels inline
\end{asy}

Then $x=-4$ and the largest $k$ is $-16$. The answer is $\boxed{-16}$.
